\documentclass{report}
\usepackage{amsmath,amssymb}
\setlength{\parindent}{0mm}
\setlength{\parskip}{1em}
\begin{document}
\begin{center}
\rule{6in}{1pt} \
{\large Patrick Greene \\
{\bf Newton's Method with Discontinuous Galerkin Projections Applied to Minimization Problems in Level Set Re-distancing and Weighted Condition Number Mesh Relaxation}}

P O Box 808 \\ L-405 \\ Livermore \\ CA 94551
\\
{\tt greene30@llnl.gov}\\
Robert Nourgaliev \\
Sam Schofield\end{center}

We demonstrate how Newton's method can be effectively combined with
discontinuous Galerkin projections for minimization problems based on
discretized quantities. Utilizing a discontinuous Galerkin projection
provides a simple method to compute the derivatives needed for the
Jacobian and Hessian since the derivatives are simply evaluated from the
prescribed basis functions. The resulting projected function and its
derivatives can then be used with Newton's method to solve the
minimization problem. Two applications of the proposed method are shown.

The first application of the method is for re-distancing in a high-order
unstructured mesh level set method. The re-distancing process requires
finding the shortest distance from a given point to the zero contour of
the level set. We introduced a new re-distancing algorithm, where this
process is formulated as a constrained minimization problem with the
level set represented as a discontinuous Galerkin projection. Static and
dynamic front evolution results will be shown, confirming that the
computed distance function maintains the same order of convergence as the
level set projection.

The second application is weighted condition number (CN) mesh relaxation.
CN mesh relaxation is used in arbitrary Lagrangian-Eulerian (ALE)
simulations to prevent mesh tangling, which can result from the
Lagrangian mesh motion. The method works by defining a condition number
functional at each vertex, which is based on the near by cell geometries,
and then iteratively minimizing the functional to find new vertex
positions. The inclusion of a weight allows the mesh relaxer to
concentrate mesh density near regions of interest. A new formulation of
the weights using a discontinuous Galerkin projection is presented. The
new method can cluster grid cells near interfaces prescribed by discrete
cell-centered values, such as an index function or volume fraction.
Sample meshes demonstrating the method�s ability to cluster cells around
dynamically evolving complex geometries are given.


This work was performed under the auspices of the U.S. Department of
Energy by Lawrence Livermore National Laboratory under Contract
DE-AC52-07NA27344. Information management release number LLNL-ABS-681162.


\end{document}
